\section{Main result}\label{sec:main}

Throughout, $(\gbar,\phibar)$ denotes the solution of the Einstein--scalar field equations
$\operatorname{Ric}_{g}=2\,d\phi\otimes d\phi$, $\Box_{g}\phi=0$ constructed by Reiterer and
Trubowitz~\cite{RT2019} on the spherically symmetric manifold $\Mhat$ of \cref{sec:setting}, with
self-similarity constant $K\approx 1.7227$ and null rescaling constant $\mu\approx 0.16831$ (both known to
$80$ digits). As in~\cite{RT2019}, we call it the \emph{Choptuik spacetime}. The discrete self-similarity
$\Theta$ acts by $\tau\mapsto\tau+2\pi$ in the coordinates $(\tau,\xi)$ of RT, rescales lengths by $e^{-K}$
and reverses the sign of $\phibar$; the full period of the background is therefore $\Theta^{2}$, that is,
$\tau\mapsto\tau+4\pi$.

A perturbation is of \emph{Floquet type with exponent $s\in\C$} if it is $e^{s\tau}$ times a profile which is
$4\pi$-periodic in $\tau$ (\cref{def:physical-mode}). Since $e^{i\tau/2}$ is itself $4\pi$-periodic, the
exponent is defined only modulo $\tfrac{i}{2}\Z$; we call the classes of $\C/\tfrac{i}{2}\Z$ \emph{Floquet
classes}, and count everything per class. In terms of the logarithmic time $T$ in which $\Theta$ is the
translation by $K$, the factor $e^{s\tau}$ is $e^{\lambda T}$ with
\begin{equation}\label{eq:lambda}
  \lambda=\frac{2\pi s}{K}.
\end{equation}

The family $L(s)$, $s\in\C$, is the linearization of the Reiterer--Trubowitz system about the background,
acting on Floquet profiles (\cref{def:L}); $C(s)$ is the corresponding linearized constraint map. Physical
Floquet modes and gauge transformations are defined in \cref{sec:physical}
(\cref{def:physical-mode,def:gauge}).

\begin{mainthm}
Let $(\gbar,\phibar)$ be the Choptuik spacetime of~\cite{RT2019}. Consider spherically symmetric linear
perturbations of Floquet type, with $L$ realized in $\Yp$ or, equivalently in the region considered, in the torus
space (\cref{sec:setting-realizations}); physical modes in the regularity class of \cref{def:physical-mode}, smooth in
$\tau$ and real analytic in $\xi$ up to and across the light cone; and gauge transformations as in
\cref{def:gauge}.
\begin{enumerate}[label=\textup{(\arabic*)},leftmargin=*]
\item \textbf{Roots of $L$.} Modulo $\tfrac{i}{2}\Z$, the roots of $L$ in the half plane $\Re s>0$, counted
  with algebraic multiplicity, are exactly four. Each of them is algebraically simple and has a real
  representative:
  \begin{center}
  \begin{tabular}{@{}lll@{}}
  \toprule
  root & enclosure & eigenvector of $L$ \\
  \midrule
  $s=\mu$ & $\mu=0.168307078963\ldots$ (exact) & $g_*=(1+Z)\omegabar$, pure gauge \\
  $s=\sB$ & $|\sB-0.0414194105|\le 3.7\cdot10^{-5}$ & violates the constraints \\
  $s=\sK$ & $|\sK-0.401024733|\le 1.72\cdot10^{-4}$ & violates the constraints \\
  $s=\sphys$ & $\sphys\in[0.73316949,\ 0.73318983]$ & satisfies the constraints \\
  \bottomrule
  \end{tabular}
  \end{center}
  On the imaginary axis, the only root of $L$ is $s\equiv 0$; it is algebraically simple, with eigenvector
  $\partial_\tau\omegabar$, the infinitesimal translation of the phase.
\item \textbf{Physical modes.} The space of physical Floquet modes with $\Re s>0$, modulo gauge, is
  one-dimensional: it consists of the modes with exponent $\sphys$, and there are no generalized modes
  \textup{(}Jordan chains\textup{)} of length $\ge2$. If the past light cone of the singularity is held fixed
  (\cref{def:conefixed}), the quotient by the smaller group of cone-preserving gauge transformations is
  two-dimensional: the additional class has exponent $\mu$ and is tangent to the one-parameter family of
  exact solutions obtained by translating the singular point, all of which are isometric to
  $(\gbar,\phibar)$. It is therefore not a second instability of the critical solution, but the effect of
  fixing the time of the singularity.
\end{enumerate}
\end{mainthm}

Here $Z=\mu^{-1}\partial_\tau+\xi\partial_\xi$. The eigenvector $g_*$ is a profile; the corresponding Floquet
mode $e^{\mu\tau}g_*$ is the Lie derivative of the background along $X_0=\partial_{u_-}+\partial_{u_+}=e^{\mu\tau}(\mu^{-1}\partial_\tau+\xi\partial_\xi)$,
the generator of translations of the singular point $(u_-,u_+)=(0,0)$ (\cref{lem:cone}).

\begin{remark}[The critical exponent]\label{rem:gamma}
By~\eqref{eq:lambda} and the value of $K$ from~\cite{RT2019}, the unique unstable physical mode grows like
$e^{\lambda T}$ with
\[
  \lambda=\frac{2\pi\sphys}{K}\in[2.674040,\ 2.674115],\qquad
  \frac{1}{\lambda}\in[0.373955,\ 0.373966].
\]
The scaling argument of critical collapse relates a single unstable mode to the mass-scaling exponent by
$\gamma=1/\lambda$~\cite{KoikeHaraAdachi1995,Gundlach1997}. Choptuik measured $\gamma\approx0.37$ in his
simulations~\cite{Choptuik1993}. Gundlach computed the unstable mode by numerical perturbation theory and found
$\lambda_1=-2.674\pm0.009$, hence $\gamma=-1/\lambda_1=0.374\pm0.001$~\cite{Gundlach1997}; the sign is that of his
time variable, in which the growing mode has negative exponent, so his $-\lambda_1$ corresponds to our
$\lambda$. Our interval lies within his error bar. This agreement is an external check, not part of the theorem: the
scaling argument concerns the nonlinear evolution, which we do not study.
\end{remark}

\begin{remark}[What is, and is not, asserted]\label{rem:scope}
\begin{enumerate}[label=(\alph*),leftmargin=*]
\item The theorem is \emph{spectral}: it counts Floquet modes of the linearized equations. It asserts
  nothing about the linearized evolution semigroup, nor about the nonlinear dynamics.
\item Only spherically symmetric perturbations are considered. Nonspherical perturbations are discussed
  in \cref{sec:discussion}.
\item Physical modes are taken in the class of \cref{def:physical-mode}: smooth in $\tau$ and real
  analytic in $\xi$, with a complex neighbourhood uniform in $\tau$. Merely smooth dependence on $\xi$
  across the light cone and at the axis is not covered (\cref{sec:discussion}).
\item Reiterer and Trubowitz prove existence of $(\gbar,\phibar)$ near their numerical data; local
  uniqueness is not formally stated in~\cite{RT2019}. The theorem is about their solution.
\item The roots of $L$ which violate the constraints are not physical: the corresponding solutions of
  the evolution part of the system do not solve the Einstein equations. The root $\mu$ is a coordinate
  effect of the moving light cone: its mode is pure gauge, and it survives only in the quotient by the smaller
  group of cone-preserving gauge transformations, which is the second count in~(2).
\end{enumerate}
\end{remark}

\begin{remark}[Neutral modes]\label{rem:neutral}
Part~(1) on the imaginary axis concerns the roots of $L$ only. At $s=0$ there are two further exact
symmetries of the equations, the global scaling $\delta g=2\gbar$ and the shift $\delta\phi=\mathrm{const}$,
which are physical perturbations with $\delta\omega=0$ and are therefore invisible to $L$. The gauge and
reconstruction lemmas of \cref{sec:physical} are proved for $\Re s>0$ only, so part~(2) has no analogue on
the axis.
\end{remark}

\begin{remark}[Dependence on computation]\label{rem:computation}
The proof is computer assisted in two places: the existence of the background with explicit bounds, which
we take from~\cite{RT2019}, and the bounds of \cref{sec:counting,sec:roots}, which are new. All of the
latter are produced in interval arithmetic (Arb~\cite{Arb}) or exact rational arithmetic, or as a posteriori
bounds on floating-point computations with rigorous rounding control. Certificates, the verification
program and the data are publicly available (\cref{sec:computation}).
\end{remark}

\subsection{Structure of the proof}\label{sec:main-structure}
The proof has three layers.
\begin{enumerate}[label=(\roman*),leftmargin=*]
\item \emph{Geometry} (\cref{sec:physical}). Every physical Floquet mode with $\Re s>0$ can be put, by a
  gauge transformation, into the form used by Reiterer and Trubowitz, and then corresponds to an element of
  $\ker L(s)\cap\ker C(s)$; conversely every such element comes from a physical mode. The only gauge modes
  with $\Re s>0$ are the multiples of $e^{\mu\tau}g_*$ (\cref{lem:gauge}). This reduces
  part~(2) to part~(1) together with the behaviour of the constraints on each root space.
\item \emph{Counting} (\cref{sec:counting}). The number of roots of $L$ in the right half plane is
  computed by an operator version of Rouch\'e's theorem with a finite-dimensional reference, after a
  rank-two (Brauer) deflation of the two gauge roots $0$ and $\mu$.
\item \emph{Identification} (\cref{sec:roots}). Each root is localized and shown to be simple by a
  Newton--Kantorovich argument, and its eigenvector is shown either to violate the constraints (by an
  explicit lower bound), or to be pure gauge, or, for $\sphys$, to satisfy them on the whole root space
  (through the injectivity of the constraint propagation operator $K$).
\end{enumerate}
